\par\Needspace{10\baselineskip}\medskip
\hypertarget{native-128-article095}{}\pdfbookmark[2]{95.}{article-095}
\begin{center}95.\end{center}
Wir kommen zur zweiten Aufgabe, der Bestimmung der Elemente aus dem gefundenen Winkel $g$. Die große Halbachse erhält man sogleich aus den Formeln 10 und $10^{\ast}$. An ihrer Stelle lassen sich auch die folgenden Formeln verwenden.
\[
\begin{aligned}
\mathrm{17)}\qquad a&=\dfrac{2mm\cos f\mathpunct{.}\BareRoot{rr'}}{yy\sin g^2}=\dfrac{kktt}{4yyrr'\cos f^2\sin g^2}\\[6pt]
\mathrm{17}^{\ast}\mathrm{)}\qquad a&=\dfrac{-2MM\cos f\mathpunct{.}\BareRoot{rr'}}{YY\sin g^2}=\dfrac{kktt}{4YYrr'\cos f^2\sin g^2}\mathpunct{\cdot}
\end{aligned}
\]
Die kleine Halbachse $b=\BareRoot{ap}$ erhält man aus Gleichung 1, die in Verbindung mit den vorhergehenden Gleichungen Folgendes ergibt
\[
\begin{aligned}
\mathrm{18)}\qquad p&=\left(\dfrac{yrr'\sin 2f}{kt}\right)^2\\[6pt]
\mathrm{18}^{\ast}\mathrm{)}\qquad p&=\left(\dfrac{Yrr'\sin 2f}{kt}\right)^2\mathpunct{\cdot}
\end{aligned}
\]
Der Flächeninhalt des elliptischen Sektors zwischen den beiden Radiusvektoren und dem Ellipsenbogen ist $=\frac12 kt\BareRoot{p}$, der des Dreiecks zwischen denselben Radiusvektoren und der Sehne $=\frac12 rr'\sin 2f$. Das Verhältnis des Sektors zum Dreieck ist daher $y:1$ oder $Y:1$. Diese Beobachtung ist von größter Bedeutung und veranschaulicht zugleich die Gleichungen 12 und $12^{\ast}$ besonders schön. Man sieht nämlich, dass in Gleichung 12 die Terme $m$, $(l+x)^{\frac12}$, $X(l+x)^{\frac32}$ und in Gleichung $12^{\ast}$ die Terme $M$, $(L-x)^{\frac12}$, $X(L-x)^{\frac32}$ jeweils dem Flächeninhalt des Sektors (zwischen den Radiusvektoren und dem Ellipsenbogen), des Dreiecks (zwischen den Radiusvektoren und der Sehne) und des Segments (zwischen dem Bogen und der Sehne) proportional sind. Denn der erste Flächeninhalt ist offensichtlich gleich der Summe oder der Differenz der beiden übrigen, je nachdem, ob $v'-v$ zwischen 0 und $180^0$ oder zwischen $180^0$ und $360^0$ liegt. Ist $v'-v$ größer als $360^0$, so muss man für jeden vollen Umlauf in dieser Bewegung zum Flächeninhalt des \SourceLeaf{129}Sektors und ebenso zu dem des Segments den Flächeninhalt der ganzen Ellipse hinzurechnen.

Da $b$ den Wert $=a\cos\varphi$ hat, folgt aus der Verbindung der Gleichungen 1, 10 und $10^{\ast}$ ferner
\[
\begin{aligned}
\mathrm{19)}\qquad \cos\varphi&=\dfrac{\sin g\operatorname{tang}f}{2(l+\sin\frac12 g^2)}\\[6pt]
\mathrm{19}^{\ast}\mathrm{)}\qquad \cos\varphi&=\dfrac{-\sin g\operatorname{tang}f}{2(L-\sin\frac12 g^2)}\mathpunct{\raisebox{.45ex}{,}}
\end{aligned}
\]
Setzt man für $l$ und $L$ ihre Werte aus Art. 89 ein, so ergibt sich daraus
\[
\mathrm{20)}\qquad \cos\varphi=\dfrac{\sin f\sin g}{1-\cos f\cos g+2\operatorname{tang}2\omega^2}\mathpunct{\cdot}
\]
Für die genaue Berechnung der Exzentrizität eignet sich diese Formel bei kleinen Werten der Exzentrizität nicht. Aus ihr lässt sich leicht die folgende, besser geeignete Formel ableiten.
\[
\mathrm{21)}\qquad \operatorname{tang}\tfrac12\varphi^2=\dfrac{\sin\frac12(f-g)^2+\operatorname{tang}2\omega^2}{\sin\frac12(f+g)^2+\operatorname{tang}2\omega^2}\mathpunct{\raisebox{.45ex}{,}}
\]
Durch Multiplikation von Zähler und Nenner mit $\cos 2\omega^2$ kann man ihr auch die folgende Form geben.
\[
\mathrm{22)}\qquad \operatorname{tang}\tfrac12\varphi^2=\dfrac{\sin\frac12(f-g)^2+\cos\frac12(f-g)^2\sin 2\omega^2}{\sin\frac12(f+g)^2+\cos\frac12(f+g)^2\sin 2\omega^2}\mathpunct{\cdot}
\]
Mit beiden Formeln (gegebenenfalls unter Verwendung von Hilfswinkeln, deren Tangenten bei der ersten Formel $\dfrac{\operatorname{tang}2\omega}{\sin\frac12(f-g)}\mathpunct{\raisebox{.45ex}{,}}$ $\dfrac{\operatorname{tang}2\omega}{\sin\frac12(f+g)}$ oder bei der zweiten Formel $\dfrac{\sin 2\omega}{\operatorname{tang}\frac12(f-g)}\mathpunct{\raisebox{.45ex}{,}}$ $\dfrac{\sin 2\omega}{\operatorname{tang}\frac12(f+g)}$ sind) lässt sich der Winkel $\varphi$ stets mit voller Genauigkeit bestimmen.

Zur Bestimmung des Winkels $G$ kann man die folgende Formel verwenden. Sie folgt unmittelbar aus der Verbindung der Gleichungen 5 und 7 mit der auf die letztere folgenden unnummerierten Gleichung.
\[
\mathrm{23)}\qquad \operatorname{tang}G=\dfrac{(r'-r)\sin g}{(r'+r)\cos g-2\cos f\mathpunct{.}\BareRoot{rr'}}\mathpunct{\raisebox{.45ex}{,}}
\]
Führt man $\omega$ ein, so erhält man daraus leicht
\[
\mathrm{24)}\qquad \operatorname{tang}G=\dfrac{\sin g\sin 2\omega}{\cos 2\omega^2\sin\frac12(f-g)\sin\frac12(f+g)+\sin 2\omega^2\cos g}\mathpunct{\cdot}
\]
Die verbleibende Mehrdeutigkeit lässt sich mithilfe von Gleichung 7 leicht auflösen. Sie zeigt, dass $G$ zwischen 0 und $180^0$ oder zwischen $180^0$ und $360^0$ zu wählen ist, je nachdem, ob der Zähler in diesen beiden Formeln positiv oder negativ ist.

\SourceLeaf{130}Aus der Verbindung von Gleichung 3 mit den folgenden Gleichungen, die sich unmittelbar aus Gleichung II in Art. 8 ergeben,
\[
\begin{aligned}
\dfrac1r-\dfrac1{r'}&=\dfrac{2e}{p}\sin f\sin F\\[6pt]
\dfrac1r+\dfrac1{r'}&=\dfrac2p+\dfrac{2e}{p}\cos f\cos F,
\end{aligned}
\]
erhält man ohne Mühe die folgende Formel.
\[
\mathrm{25)}\qquad \operatorname{tang}F=\dfrac{(r'-r)\sin f}{2\cos g\mathpunct{.}\BareRoot{rr'}-(r'+r)\cos f}\mathpunct{\raisebox{.45ex}{,}}
\]
Führt man in ihr den Winkel $\grave{\omega}$ ein, so ergibt sich
\[
\mathrm{26)}\qquad \operatorname{tang}F=\dfrac{\sin f\sin 2\omega}{\cos 2\omega^2\sin\frac12(f-g)\sin\frac12(f+g)-\sin 2\omega^2\cos f}\mathpunct{\cdot}
\]
Die Mehrdeutigkeit wird hier ebenso wie zuvor aufgehoben. — Sobald die Winkel $F$ und $G$ gefunden sind, erhält man $v=F-f$, $v'=F+f$, womit die Lage des Perihels bekannt \raisebox{.45ex}{.}ist. Ferner erhält man $E=G-g$, $E'=G+g$. Schließlich hat die mittlere Bewegung während der Zeit $t$ den Wert $=\dfrac{kt}{a^{\frac32}}=2g-2e\cos G\sin g$. Die Übereinstimmung dieser Ausdrücke dient zur Kontrolle der Rechnung. Die Epoche der mittleren Anomalie für den Zeitpunkt in der Mitte zwischen den beiden gegebenen Zeitpunkten ist $G-e\sin G\cos g$ und kann nach Belieben auf jeden anderen Zeitpunkt übertragen werden. Noch etwas bequemer ist es, die mittleren Anomalien für die beiden gegebenen Zeitpunkte mit den Formeln $E-e\sin E$ und $E'-e\sin E'$ zu berechnen und ihre Differenz zur Kontrolle der Rechnung mit $\dfrac{kt}{a^{\frac32}}$ zu vergleichen.
